\begin{explanationbox}
\textbf{先看具体曲线：为什么要加二次项？}
\begin{center}
\begin{tikzpicture}[font=\small]
\draw[black!30,rounded corners=2pt] (0,0) rectangle (14.4,8.2);
\node[anchor=west,font=\small\bfseries] at (0.5,7.72) {同一点出发：二次项补上指数曲线的弯曲};
\begin{axis}[width=11.2cm,height=6.1cm,at={(2.05cm,1.25cm)},anchor=south west,
axis lines=left,xlabel={$x$},ylabel={$y$},tick label style={font=\small},label style={font=\small},
legend style={font=\small,draw=none,fill=white,at={(0.03,0.97)},anchor=north west},
legend cell align=left,unbounded coords=jump,xmin=-0.8,xmax=0.8,ymin=0,ymax=2.5]
\addplot[blue,very thick] table[x=x,y=exp] {assets/tikz/F004_samples.dat};
\addlegendentry{$e^x$}
\addplot[orange!85!black,thick,dashed] table[x=x,y=pe] {assets/tikz/F004_samples.dat};
\addlegendentry{$1+x+x^2/2$}
\addplot[purple,thick,dashdotted] table[x=x,y=line] {assets/tikz/F004_samples.dat};
\addlegendentry{$1+x$}

\end{axis}

\end{tikzpicture}
\end{center}

直线 $1+x$ 与 $e^x$ 在 $(0,1)$ 处同高、同斜率；加上 $x^2/2$，还匹配了二阶导数。
所谓“贴合”，指这些导数在展开点相同，并有余项控制；不能据此宣称它在任意区间内误差最小。

\textbf{同样的办法，用在不同形状上。}
\begin{center}
\begin{tikzpicture}[font=\small]
\draw[black!30,rounded corners=2pt] (0,0) rectangle (14.4,8.2);
\node[anchor=west,font=\small\bfseries] at (0.5,7.72) {同一尺度：正弦过原点，余弦从 $(0,1)$ 出发};
\begin{axis}[width=5.2cm,height=5.3cm,at={(1.85cm,1.25cm)},anchor=south west,
axis lines=left,xlabel={$x$},ylabel={$y$},xmin=-1.2,xmax=1.2,ymin=-1.1,ymax=1.25,
tick label style={font=\small},label style={font=\small},
legend style={font=\small,draw=none,fill=none,at={(0.5,1.15)},anchor=south},legend cell align=left]
\addplot[blue,very thick] table[x=x,y=sin] {assets/tikz/F004_samples.dat};
\addlegendentry{$\sin x$}
\addplot[orange!85!black,thick,dashed] table[x=x,y=ps] {assets/tikz/F004_samples.dat};
\addlegendentry{$x-x^3/6+x^5/120$}
\end{axis}
\begin{axis}[width=5.2cm,height=5.3cm,at={(8.25cm,1.25cm)},anchor=south west,
axis lines=left,xlabel={$x$},ylabel={$y$},xmin=-1.2,xmax=1.2,ymin=-1.1,ymax=1.25,
tick label style={font=\small},label style={font=\small},
legend style={font=\small,draw=none,fill=none,at={(0.5,1.15)},anchor=south},legend cell align=left]
\addplot[blue,very thick] table[x=x,y=cos] {assets/tikz/F004_samples.dat};
\addlegendentry{$\cos x$}
\addplot[orange!85!black,thick,dashed] table[x=x,y=pc] {assets/tikz/F004_samples.dat};
\addlegendentry{$1-x^2/2+x^4/24$}
\end{axis}

\end{tikzpicture}
\end{center}

正弦为奇函数，余弦为偶函数，展开中相应缺少偶次项或奇次项。
图中两类函数使用相同坐标范围；曲线几乎重合时，后面的误差图更能分辨精度。

\begin{center}
\begin{tikzpicture}[font=\small]
\draw[black!30,rounded corners=2pt] (0,0) rectangle (14.4,8.2);
\node[anchor=west,font=\small\bfseries] at (0.5,7.72) {正切：三次修正为正，仍只在零点附近贴合};
\begin{axis}[width=11.2cm,height=6.1cm,at={(2.05cm,1.25cm)},anchor=south west,
axis lines=left,xlabel={$x$},ylabel={$y$},tick label style={font=\small},label style={font=\small},
legend style={font=\small,draw=none,fill=white,at={(0.03,0.97)},anchor=north west},
legend cell align=left,unbounded coords=jump,xmin=-1.2,xmax=1.2,ymin=-3,ymax=3,xtick={-1,-0.5,0,0.5,1},restrict x to domain=-1.2:1.2]
\addplot[blue,very thick] table[x=x,y=tan] {assets/tikz/F004_samples.dat};
\addlegendentry{$\tan x$}
\addplot[orange!85!black,thick,dashed] table[x=x,y=pt] {assets/tikz/F004_samples.dat};
\addlegendentry{$x+x^3/3$}

\end{axis}

\end{tikzpicture}
\end{center}

正切比 $x$ 增长得快，三次项应为正。它接近 $\pm\pi/2$ 时无界，有限多项式不会复制这种竖直渐近线。

\begin{center}
\begin{tikzpicture}[font=\small]
\draw[black!30,rounded corners=2pt] (0,0) rectangle (14.4,8.2);
\node[anchor=west,font=\small\bfseries] at (0.5,7.72) {对数：靠近边界 $x=-1$ 时，三次近似明显偏离};
\begin{axis}[width=11.2cm,height=6.1cm,at={(2.05cm,1.25cm)},anchor=south west,
axis lines=left,xlabel={$x$},ylabel={$y$},tick label style={font=\small},label style={font=\small},
legend style={font=\small,draw=none,fill=white,at={(0.03,0.97)},anchor=north west},
legend cell align=left,unbounded coords=jump,xmin=-1.05,xmax=0.8,ymin=-4.2,ymax=1.1,legend style={at={(0.97,0.03)},anchor=south east}]
\addplot[blue,very thick] table[x=x,y=log] {assets/tikz/F004_samples.dat};
\addlegendentry{$\ln(1+x)$}
\addplot[orange!85!black,thick,dashed] table[x=x,y=pl] {assets/tikz/F004_samples.dat};
\addlegendentry{$x-x^2/2+x^3/3$}
\draw[black!50,densely dotted] (axis cs:-1,-4.1)--(axis cs:-1,1);
\node[anchor=west] at (axis cs:-0.93,-3.7) {$x\to-1^+$};
\end{axis}

\end{tikzpicture}
\end{center}

零点附近的三次多项式不能复制 $x\to-1^+$ 时对数趋于 $-\infty$ 的行为。
所以“$|x|$ 多小才够用”必须由精度要求和误差界决定，没有适用于所有函数的固定阈值。
\end{explanationbox}
